%% backmatter/glossaire.tex — Glossaire du livre « Claude Code en pratique »
%% Commande sémantique : \entreeglossaire{terme}{définition}
%% définie dans theme-technique.sty

\chapter*{Glossaire}
\addcontentsline{toc}{chapter}{Glossaire}
\markboth{Glossaire}{Glossaire}

\entreeglossaire{Agent}{Programme fondé sur un modèle de langage qui agit en
boucle : il perçoit son environnement (fichiers, commandes, messages), décide
d'une action, l'exécute, observe le résultat, puis recommence jusqu'à atteindre
son but.}

\entreeglossaire{Claude Code}{Agent de développement en ligne de commande
d'Anthropic. Opère directement dans le dépôt : lit le code, propose et applique
des modifications, lance des commandes, et demande l'approbation du développeur
selon les règles configurées.}

\entreeglossaire{Contexte}{L'ensemble des informations qu'un agent a
effectivement sous les yeux à un instant donné : instructions, historique de la
session, fichiers lus, résultats de commandes. Tout ce qui n'y figure pas
n'existe pas pour l'agent.}

\entreeglossaire{Fenêtre de contexte}{La quantité maximale d'information qu'un
agent peut garder active simultanément. Quand elle approche de la saturation,
l'agent commence à oublier les éléments les plus anciens et ses réponses se
dégradent.}

\entreeglossaire{Garde-fou}{Règle ou mécanisme qui limite la portée des actions
d'un agent pour contenir les erreurs et les rendre réversibles. Exemples :
exiger une confirmation avant toute modification, restreindre les commandes
autorisées, travailler sur une branche dédiée.}

\entreeglossaire{Hooks}{Scripts shell déclenchés automatiquement par Claude Code
avant ou après certains événements (soumission d'un message, arrêt de l'agent,
etc.). Permettent d'automatiser des vérifications ou des actions récurrentes.}

\entreeglossaire{MCP (Model Context Protocol)}{Protocole standardisé qui permet
à Claude Code de se connecter à des outils et sources de données externes
(bases de données, APIs, services) via des serveurs MCP dédiés.}

\entreeglossaire{Mémoire projet}{Fichier d'instructions persistantes (typiquement
\code{CLAUDE.md}) rangé dans le dépôt et chargé automatiquement par l'agent à
chaque session. Contient les conventions du projet, les commandes de test, les
pièges connus.}

\entreeglossaire{Orchestration}{Coordination de plusieurs agents ou sous-agents
travaillant en parallèle sur des tâches complémentaires, pilotée par un agent
coordinateur.}

\entreeglossaire{Permissions}{Configuration qui détermine quelles actions Claude
Code peut exécuter sans demander de confirmation (lecture de fichiers, exécution
de commandes shell, modification de fichiers, etc.).}

\entreeglossaire{Sous-agent (subagent)}{Instance d'agent lancée par un agent
coordinateur pour traiter une sous-tâche en parallèle. Chaque sous-agent opère
de façon indépendante dans son propre contexte.}

\entreeglossaire{Token}{Unité de base du traitement par un modèle de langage.
Correspond approximativement à trois quarts d'un mot en français. La fenêtre de
contexte et les coûts d'API sont mesurés en tokens.}
